\documentclass[letterpaper,11pt]{article}

% Shared LaTeX preamble for CV and Cover Letter

\usepackage{iftex}

\ifPDFTeX
  \usepackage[T2A]{fontenc}
  \usepackage[utf8]{inputenc}
\else
  \usepackage{fontspec}
\fi

\usepackage[english, russian]{babel}

\usepackage{latexsym}
\usepackage[empty]{fullpage}
\usepackage{titlesec}
\usepackage{marvosym}
\usepackage[usenames,dvipsnames]{color}
\usepackage{verbatim}
\usepackage{enumitem}
\usepackage[hidelinks]{hyperref}
\usepackage{fancyhdr}
\usepackage{tabularx}
\usepackage{graphicx}

\input{glyphtounicode}
\pdfgentounicode=1

\pagestyle{fancy}
\fancyhf{} % clear all header and footer fields
\fancyfoot{}
\renewcommand{\headrulewidth}{0pt}
\renewcommand{\footrulewidth}{0pt}

\urlstyle{same}
\raggedbottom
\raggedright

\newcommand{\username}{alexo375}
\newcommand{\github}{https://github.com/\username}
\newcommand{\cpp}{C\texttt{++}}


% Adjust margins
\addtolength{\oddsidemargin}{-0.5in}
\addtolength{\evensidemargin}{-0.5in}
\addtolength{\textwidth}{1in}
\addtolength{\topmargin}{-.7in} % было "-.5in". Подняли текст чуть выше
\addtolength{\textheight}{1.4in} % было "1.0in". Увеличили рабочую область страницы

\setlength{\tabcolsep}{0in}

% Sections formatting
\titleformat{\section}{
  \vspace{-4pt}\scshape\raggedright\large
}{}{0em}{}[\color{black}\titlerule \vspace{-5pt}]

%-----------CUSTOM-COMMANDS-----------
\newcommand{\resumeItem}[1]{
  \item\small{
    {#1 \vspace{-2pt}}
  }
}

\newcommand{\resumeSubheading}[4]{
  \vspace{-2pt}\item
    \begin{tabular*}{0.97\textwidth}[t]{l@{\extracolsep{\fill}}r}
      \textbf{#1} & #2 \\
      \textit{\small#3} & \textit{\small #4} \\
    \end{tabular*}\vspace{-7pt}
}

\newcommand{\resumeProjectHeading}[2]{
    \item
    \begin{tabular*}{0.97\textwidth}{l@{\extracolsep{\fill}}r}
      \small#1 & #2 \\
    \end{tabular*}\vspace{-7pt}
}

\newcommand{\resumeSubHeadingListStart}{\begin{itemize}[leftmargin=0.15in, label={}]}
\newcommand{\resumeSubHeadingListEnd}{\end{itemize}}
\newcommand{\resumeItemListStart}{\begin{itemize}}
\newcommand{\resumeItemListEnd}{\end{itemize}\vspace{-5pt}}

\begin{document}

%----------ЗАГОЛОВОК----------
\begin{center}
  \begin{minipage}[b]{0.80\textwidth}
      \textbf{\Huge \scshape Горобченко Александр} \\ \vspace{4pt}
      \large Junior / Intern ML Engineer \\ \vspace{4pt}
      \small \href{tel:+79953755007}{+7 (995) 375-50-07} $|$ 
      \href{mailto:alexo375@yandex.ru}{\underline{alexo375@yandex.ru}} $|$ 
      \href{https://t.me/alexo375}{\underline{t.me/alexo375}} $|$
      \href{https://github.com/almost-wizard}{\underline{github.com/almost-wizard}} \\ \vspace{2pt}
      \small Рязань / Удаленная работа
  \end{minipage}%
  \begin{minipage}[b]{0.18\textwidth}
      \raggedleft
      \includegraphics[width=4cm, keepaspectratio]{me.png} 
  \end{minipage}
\end{center}

%----------О-СЕБЕ----------
\section{О Себе}
 \begin{itemize}[leftmargin=0.15in, label={}]
    \small{\item{
     ML Engineer с опытом разработки алгоритмов CV и построения ML-пайплайнов.
     Разрабатываю модели для анализа изображений, проектирую микросервисную архитектуру для ML-решений.
     Имею опыт обработки спутниковых изображений на Python и C++.
     Ориентирован на создание масштабируемых ML систем.
    }} 
 \end{itemize}

%-----------НАВЫКИ-----------
\section{Технические навыки}
 \begin{itemize}[leftmargin=0.15in, label={}]
    \small{\item{
     \textbf{Python, ML \& CV}{: PyTorch, CatBoost, XGBoost, LightGBM, Scikit-Learn, OpenCV, Pandas, NumPy, SciPy, GDAL} \\
     \textbf{C++}{: CMake, LibTorch, Boost, GTest, Valgrind, Conan} \\
     \textbf{Data Engineering}{: PostgreSQL, Greenplum, ClickHouse, OracleDB, MongoDB, Amazon S3, Airflow, Spark, Kafka} \\
     \textbf{DevOps \& CI/CD}{: Docker, Git, Github Actions, Linux (NixOS, Ubuntu)}
    }}
 \end{itemize}

%-----------ОПЫТ-----------
\section{Опыт работы}
  \resumeSubHeadingListStart
    \resumeSubheading
      {НИИ Фотон, ФГБОУ ВО РГРТУ}{Рязань}
      {R\&D Инженер (C++ / Python / ML)}{Июль 2023 - Июль 2026}
      \resumeItemListStart
        \resumeItem{Разработал классические и ML алгоритмы обнаружения изменений на спутниковых
        снимках площадной съемки высокого разрешения. С нуля подготовил и аннотировал обучающую выборку для CNN-модели.}
        \resumeItem{Опубликовал 3 научные статьи в рамках конференций
        \href{https://elibrary.ru/item.asp?id=82496414}{\underline{СТНО-2025}},
        \href{https://elibrary.ru/item.asp?id=87261582}{\underline{НИТ-2025}} и
        СТНО-2026.}
      \resumeItemListEnd
  \resumeSubHeadingListEnd

%-----------ПРОЕКТЫ-----------
\section{Проекты}
    \resumeSubHeadingListStart
      \resumeProjectHeading
          {\textbf{Change Detector} $|$ \emph{C++, Python, PyTorch, LibTorch}}{Август 2025 - Июнь 2026}
          \resumeItemListStart
            \resumeItem{Реализовал детерминированные алгоритмы для обнаружения изменений на спутниковых снимках.
            Оптимизировал их работу, реализовав подбор гиперпараметров биоинспирированными алгоритмами.}
            \resumeItem{Обучил CNN (UNet++, ResNet-34) для сегментации изменений на снимках,
            выполнив ручную разметку кастомного датасета (ключевые метрики: $F1 \approx 0.87$, $IoU \approx 0.76$)}
          \resumeItemListEnd
          
      \resumeProjectHeading
          {\href{https://github.com/almost-wizard/acc-2025-validation-pipeline}{\textbf{\underline{ML Image Validator}}} $|$ \emph{Python, FastAPI, ML Pipeline}}{Октябрь 2025 - Декабрь 2025}
          \resumeItemListStart
            \resumeItem{Реализовал микросервисный пайплайн для ML-валидации изображений в команде из 3 человек.}
            \resumeItem{Создал на FastAPI микросервис, обеспечивающий стабильную проверку качества данных.}
          \resumeItemListEnd

      \resumeProjectHeading
          {\href{https://github.com/almost-wizard/native-compiler}{\textbf{\underline{Native-compiler}}} $|$ \emph{C++}}{Октябрь 2025 - Декабрь 2025}
          \resumeItemListStart
            \resumeItem{Разработал консольную утилиту, демонстрирующую работу основных фаз компилятора.}
            \resumeItem{Реализовал лексический, семантический, синтаксический анализ; генератор 3-адресного кода, постфиксной формы.}
          \resumeItemListEnd
    \resumeSubHeadingListEnd

%-----------ОБРАЗОВАНИЕ-----------
\section{Образование}
  \resumeSubHeadingListStart
    \resumeSubheading
      {Центральный Университет}{Москва}
      {Математика и компьютерные науки, трек ML-online $|$ GPA: 9.25 $|$ Магистратура}{2025 - Настоящее время}
    \resumeItemListStart
      \resumeItem{Релевантные курсы: ML, Data Engineering, Основы статистики}
      \resumeItem{Выиграл грант, покрывающий 50\% стоимости обучения}
    \resumeItemListEnd
    \resumeSubheading
      {РГРТУ имени В.Ф. Уткина}{Рязань}
      {Программная инженерия $|$ GPA: 4.80 $|$ Бакалавриат}{2022 - 2026}
    \resumeItemListStart
      \resumeItem{Релевантные курсы: Высшая математика, алгоритмы, низкоуровневое программирование}
      \resumeItem{Бюджетная форма обучения}
    \resumeItemListEnd
  \resumeSubHeadingListEnd
  
%-----------Дополнительно-----------
\section{Дополнительная информация}
 \begin{itemize}[leftmargin=0.15in, label={}]
    \small{\item{
     \textbf{Языки}{: Английский (B2)}
    }}
 \end{itemize}


%-------------------------------------------
\end{document}
